% Chapter 3: starting and stopping MiniDAQ's pipeline. Consumers start first; only the source
% gets a stop token; a stop ends the input, and every later stage drains and ends on its own.
\begin{tikzpicture}[x=1cm,y=1cm]
  \foreach \n/\x/\o in {source/0/5, decode/2.6/4, sort/5.2/3, pair/7.8/2, write/10.4/1} {
    \node[box,minimum width=1.9cm] (\n) at (\x,0) {\n};
    \node[note,text=d2dink,anchor=south] at (\x,0.45) {starts \o.};}
  \foreach \a/\b in {source/decode,decode/sort,sort/pair,pair/write} {\draw[flow] (\a) -- (\b);}
  \node[hot=d2dcat3,minimum width=1.6cm,font=\sffamily\scriptsize\bfseries] (st) at (0,-1.6) {stop token};
  \draw[->,d2dcat3,line width=1.1pt] (st) -- (source);
  \draw[d2dcat2,line width=1pt,->] (1.2,-0.75) -- (11.4,-0.75);
  \node[lab,anchor=north west,text=d2dink,align=left] at (1.3,-0.85) {stop: the source ends its input; each stage drains what it accepted, sees its input\\closed, and returns. No data already accepted is discarded. Each \texttt{std::jthread}'s\\destructor then requests stop and joins: no \texttt{std::terminate} on an early return.};
\end{tikzpicture}
